\documentclass[10pt,a4paper]{article}
\usepackage[T1]{fontenc}\usepackage[utf8]{inputenc}\usepackage[english,italian]{babel}
\usepackage{lmodern,microtype}\usepackage[a4paper,margin=1.85cm,headheight=15pt]{geometry}
\usepackage{enumitem,tabularx,array,longtable,booktabs,pdflscape,xcolor,hyperref,xurl,fancyhdr,titlesec,framed,listings}
\definecolor{ddtadark}{RGB}{28,38,52}\definecolor{ddtablue}{RGB}{42,88,135}\definecolor{ddtagray}{RGB}{244,246,248}\definecolor{ddtagreen}{RGB}{48,111,78}\definecolor{ddtaorange}{RGB}{148,91,25}
\hypersetup{colorlinks=true,linkcolor=ddtablue,urlcolor=ddtablue}
\pagestyle{fancy}\fancyhf{}\lhead{DDTA - BA0-R}\rhead{REVIEW-ONLY R2}\cfoot{\thepage}
\titleformat{\section}{\large\bfseries\color{ddtadark}}{\thesection}{.6em}{}\titleformat{\subsection}{\normalsize\bfseries\color{ddtadark}}{\thesubsection}{.6em}{}
\setlist[itemize]{leftmargin=1.45em,itemsep=.16em,topsep=.2em}\setlist[enumerate]{leftmargin=1.7em,itemsep=.16em,topsep=.2em}
\setlength{\parindent}{0pt}\setlength{\parskip}{.32em}\setlength{\emergencystretch}{2em}
\lstset{basicstyle=\ttfamily\small,frame=single,breaklines=true,backgroundcolor=\color{ddtagray}}
\newenvironment{challenge}[1]{\def\FrameCommand{\fcolorbox{ddtaorange}{ddtagray}}\MakeFramed{\advance\hsize-2\width\FrameRestore}\noindent\textbf{#1}\par\smallskip}{\endMakeFramed}
\newenvironment{gate}[1]{\def\FrameCommand{\fcolorbox{ddtagreen}{ddtagray}}\MakeFramed{\advance\hsize-2\width\FrameRestore}\noindent\textbf{#1}\par\smallskip}{\endMakeFramed}

\hypersetup{pdftitle={BA0-R systems-modeling prior-art synthesis Revision 2}}
\begin{document}
\begin{center}{\LARGE\bfseries BA0-R systems-modeling prior-art synthesis}\par{\large Revision 2 - REVIEW-ONLY}\par\smallskip
{\small Baseline: \texttt{89e5486e02d07ff5b97082f73a2e22ac1b4319ae}}\end{center}
\begin{challenge}{Status}
Includes BA0-R9 OPM alternative-kernel challenge. No external concept is promoted to the DDTA BAE metamodel by this document.
\end{challenge}
\section{Corpus and access boundary}
The review combines SRC-0034 and candidate SRC-0039--SRC-0047: NASA modeling guidance, SysML v2, KerML, NASA systems-engineering guidance, ISO/IEC/IEEE 42010, AADL, ArchiMate, Estefan's MBSE survey, and the reflective Object-Process Methodology (OPM) metamodel.

ISO/IEC/IEEE 42010:2022 and ArchiMate 3.2 remain \textbf{ACCESS-LIMITED}. OPM SRC-0047 uses an author-uploaded full-text rendering of the 2005 reflective metamodel; current OPM standard status is corroborated only through official ISO 19450:2024 catalogue metadata.
\section{What OPM changes}
The first reading ring showed recurring semantic responsibilities. OPM supplies the missing counterexample to the stronger inference that those responsibilities determine one necessary class hierarchy:
\begin{lstlisting}
Element
  +-- Entity
  |     +-- Thing
  |     |     +-- Object
  |     |     `-- Process
  |     `-- State [owned by Object]
  `-- Link
        +-- StructuralLink
        `-- ProceduralLink
\end{lstlisting}
System/environment, physical/informational and scope distinctions can be attributes/classifications. Behavior and flow can be carried by typed procedural links. Multiple diagrams refine or abstract one model. Therefore DDTA must separate \textbf{semantic responsibility} from \textbf{representation mechanism}.
\clearpage
\begin{landscape}
\section{Cross-source semantic-responsibility matrix}
\scriptsize
\begin{longtable}{p{3.0cm}p{6.0cm}p{5.0cm}p{7.0cm}}
\toprule Responsibility & NASA/SysML/KerML/AADL convergence & OPM counter-model & DDTA implication\\\midrule\endhead
Structure & explicit structural/part/component concepts & Object + structural links & preserve structure; exact root type OPEN\\
Behavior & function/action/behavior explicit & Process is a first-class Thing & strongly required semantic distinction\\
Information/resource & item/data/resource-flow concepts & informational essence can classify elements & separate information identity not universally forced\\
Connection & connector/connection/link concepts & Link is first-class & relation identity strongly evidenced\\
Interface/endpoint & often explicit & no universal interface root forced & first-class Interface OPEN\\
Flow/transfer & often explicit with payload & procedural links carry flow & Flow root WEAKENED\\
State/mode & often explicit & State explicit but owned by Object & State root WEAKENED; semantics remain strong\\
Boundary/environment & context/boundary/environment explicit & systemic/environmental affiliation on any Element & Boundary root WEAKENED; scope distinction strong\\
Property/constraint & explicit & attributes, characterization, timing and cardinality & likely property mechanism, not necessarily BAE\\
View/projection & views/renderings/model kinds & one model, multiple OPDs + textual modality & canonical model != view reinforced\\
Analysis & AnalysisCase / analysis-oriented models & analysis/simulation historically supported & novelty cannot be generic analysis-on-model\\
\bottomrule
\end{longtable}
\end{landscape}
\section{Hypothesis status changes}
\begin{tabularx}{\linewidth}{>{\bfseries}p{4.4cm}X}
Actor root & REJECTED AS GENERIC ROOT. Actor may survive as a contextual role/specialization.\\
Asset root & REJECTED AS GENERIC ROOT. Asset is more plausibly a security/relevance interpretation.\\
Interaction edge & REJECTED AS SUFFICIENT GENERIC RELATION. It collapses connection, interface, flow, action and procedural semantics.\\
Channel root & FURTHER WEAKENED. Communication path can be connector, bus/component, or physical Link.\\
Store root & OPEN / WEAKENED. Persistence exists but does not force a universal root.\\
Contract root & OPEN / WEAKENED. Contract-like semantics may be distributed over interfaces, schemas, constraints and requirements.\\
Boundary root & OPEN / WEAKENED. The distinction is necessary; a Boundary object is not.\\
State root & OPEN / WEAKENED. State may be owned by another element.\\
Flow root & OPEN / WEAKENED. Flow can be typed relation semantics.\\
InformationObject root & OPEN / WEAKENED AS UNIVERSAL ROOT. Stable payload identity remains a DDTA corpus question.\\
\end{tabularx}
\section{Strong semantic obligations for BA0}
A DDTA analyzable representation must be capable of preserving, without semantic collapse:
\begin{enumerate}
\item identifiable system referents;
\item behavior/transformation;
\item typed relations/connections;
\item information/resource semantics where relevant;
\item state/mode/condition semantics where relevant;
\item system/environment and responsibility/scope distinctions;
\item properties/constraints relevant to analysis;
\item abstraction/refinement and multiple projections/views from canonical knowledge.
\end{enumerate}
The mechanism remains OPEN: entity, relationship, role, property/classification, owned sub-state, or derived projection.
\section{Canonical model versus view}
NASA/SysML already supported projection. OPM adds semantically equivalent graphical/textual modalities and refinement/abstraction diagrams from one model. The working direction is therefore:
\begin{lstlisting}
canonical analyzable knowledge
    -> select / abstract / refine / project / render
    -> one or more views
\end{lstlisting}
A generated diagram or method-specific view is not semantic authority.
\section{Boundary result after OPM}
\begin{lstlisting}
system / responsibility scope distinction  STRONGLY EVIDENCED
Boundary as first-class BAE                 OPEN / WEAKENED
TrustBoundary                               OPEN; likely security overlay
\end{lstlisting}
\section{Novelty challenge}
SysML already has AnalysisCase; AADL/NASA support analysis; OPM supports model analysis/simulation and transformations. DDTA must not claim novelty for having a generic model, analysis envelope, view, or transformation. The candidate contribution remains the governed end-to-end boundary:
\begin{lstlisting}
governed documentation
 -> reviewed methodology-neutral analyzable representation
 -> replaceable method-owned security-analysis overlay
 -> reviewed normalized findings
 -> accepted governed-document / requirement change
 -> preserved analysis-to-change provenance
 -> change-aware re-analysis
\end{lstlisting}
No reviewed systems-modeling source establishes that entire loop; this is a scoped comparative inference, not proof of novelty.
\section{Chapter regression}
No Chapter 3 or Chapter 4 reopen trigger is produced by BA0-R9. OPM informs the downstream analysis layer and strengthens the requirement that first-class BAE identity must be justified. ISO 42010 full text, when supplied, should be integrated as a regression source for view/viewpoint/architecture-description claims.
\section{BA0 working hypothesis - still OPEN}
\begin{gate}{Working hypothesis}
Base Analysis is a canonical, methodology-neutral analyzable representation of governed project/system knowledge, grounded in governed documentation and explicit reviewed analytical additions. It must preserve analysis-relevant semantic distinctions without committing DDTA to one external modeling-language ontology, and support multiple consistent projections/views from one accepted knowledge base.
\end{gate}
Non-goals: second narrative hierarchy, SysML subset, OPM subset, DFD, threat-model core, raw LLM/NLP extraction, tool-native schema authority.
\section{What BA1 must test}
\begin{lstlisting}
A. typed-entity-rich kernel
   (SysML/NASA-like granularity)

B. compact entity + typed-relation + properties kernel
   (OPM-like pressure case)

C. DDTA hybrid
   first-class identity only where independent lifecycle,
   reference, provenance or analysis attachment requires it
\end{lstlisting}
The winning representation is the smallest one that preserves required DDTA semantics without ambiguity, semantic loss, generic escape relations, or method-specific contamination.
\section{Review gate}
BA0-R remains REVIEW-ONLY until user review. ISO 42010 full text may be integrated when available. No BAE metaclass is accepted by this review.
\end{document}
